% =====================================================================
%  Chapitre 2 - Etude et analyse globale du projet
% =====================================================================
\chapter{Etude et analyse globale du projet}
\label{chap:analyse}

\section*{Introduction}
\addcontentsline{toc}{section}{Introduction}

Le chapitre precedent a etabli le besoin. Celui-ci le traduit en
exigences. Il identifie d'abord les acteurs qui interagissent avec la
plateforme, puis recense les besoins fonctionnels domaine par domaine et
les besoins non fonctionnels qui conditionnent la qualite du systeme. Ces
exigences sont ensuite synthetisees par un diagramme de cas d'utilisation
global. Le chapitre se clot sur le pilotage du projet : composition de
l'equipe, carnet de produit et decoupage du travail en six sprints.

% =====================================================================
\section{Identification des acteurs}
\label{sec:acteurs}

Un acteur est une entite exterieure au systeme qui interagit avec lui.
L'analyse des habilitations reellement portees par la plateforme fait
apparaitre six acteurs humains et deux systemes tiers.

\subsection{Les acteurs principaux}

\begin{table}[H]
\centering
\caption{Acteurs principaux de la plateforme}
\label{tab:acteurs}
\small
\begin{tabularx}{\textwidth}{@{}p{2.4cm}X@{}}
\toprule
\textbf{Acteur} & \textbf{Role et perimetre d'action} \\ \midrule
Utilisateur & Acteur generique dont derivent tous les autres : compte,
  profil, second facteur, consultations non restreintes, messagerie \\
Client & Donneur d'ordre, entreprise ou particulier. Depose une demande,
  recoit un devis, suit l'acheminement, consulte ses factures, ouvre des
  tickets de support \\
Employe & Acteur \essentiel{abstrait} regroupant le personnel interne.
  Il n'agit pas : il factorise ce que partagent les roles internes \\
Chauffeur & Conduit les vehicules. Connexion par code PIN, consultation
  de ses missions, transmission de sa position \\
Technicien & Intervient sur les vehicules : interventions, consommation
  de pieces, avancement des ordres de travail, documentation \\
Responsable & Acteur \essentiel{abstrait} couvrant la combinaison
  Administrateur ou Gestionnaire, la plus frequente du code. Referentiels,
  devis, missions, maintenance, facturation \\
Administrateur & Droits exclusifs : approbation et rejet des comptes,
  habilitations, archivage, purge, suppressions definitives \\
Gestionnaire & Exploitation courante, sans droit de suppression ni
  d'administration des comptes \\
\bottomrule
\end{tabularx}
\end{table}

\subsection{La hierarchie des acteurs}

Ces acteurs ne sont pas independants : ils forment une hierarchie de
generalisation calquee sur l'heritage reel du modele de donnees, ou la
table des utilisateurs porte un discriminant distinguant les clients des
employes.

\begin{definitionbox}{Pourquoi une hierarchie d'acteurs}
Sans generalisation, un cas d'utilisation ouvert a tous devrait etre
relie separement au Client, au Chauffeur, au Technicien, au Gestionnaire
et a l'Administrateur, soit cinq associations pour une seule capacite.
La generalisation permet de ne le relier qu'a l'acteur \exemple{le plus
general qui le possede} : les acteurs specialises en heritent. Le
diagramme gagne en lisibilite sans rien perdre en precision.
\end{definitionbox}

Deux acteurs de cette hierarchie sont \essentiel{abstraits} : Employe et
Responsable ne correspondent a aucun role stocke en base. Ils
n'existent que pour factoriser. \essentiel{Responsable} merite une
explication : il traduit la combinaison d'habilitations
\exemple{Administrateur ou Gestionnaire}, qui revient dans la majorite
des controles d'acces de la plateforme. Le nommer evite de dupliquer
chaque association.

% =====================================================================
\section{Besoins fonctionnels}
\label{sec:besoins-fonctionnels}

Les besoins fonctionnels sont presentes par domaine metier. Ce decoupage
n'est pas seulement un procede de redaction : c'est celui qui structure
la plateforme et, on le verra en section~\ref{sec:sprints}, celui qui
determine le decoupage en sprints.

\begin{table}[H]
\centering
\caption{Besoins fonctionnels par domaine metier}
\label{tab:besoins-fonctionnels}
\small
\begin{tabularx}{\textwidth}{@{}p{3.1cm}X@{}}
\toprule
\textbf{Domaine} & \textbf{Besoins recenses} \\ \midrule

Identites, acces et securite &
  S'inscrire et confirmer son adresse electronique ; se connecter, avec
  validation d'un second facteur lorsqu'il est actif ; se connecter par
  code PIN depuis l'application chauffeur ; gerer son profil et son mot
  de passe ; activer ou desactiver son second facteur ; statuer sur les
  demandes d'inscription ; gerer les habilitations ; archiver puis
  restaurer un compte ; consulter et effacer l'historique des
  connexions \\

Referentiel de flotte et d'equipements &
  Consulter le parc et rechercher un vehicule disponible ; tenir la fiche
  d'un vehicule ; gerer l'inventaire des equipements, leur affectation et
  leur etat ; enregistrer les pieces administratives et les assurances et
  surveiller leurs echeances ; affecter et restituer un equipement ;
  realiser une inspection ; declarer les pleins et suivre la
  consommation ; sortir un bien de l'inventaire \\

Demande, devis, mission et suivi &
  Deposer une demande depuis le portail client avec estimation
  tarifaire ; etablir un devis en retenant, modifiant ou fixant le prix ;
  accepter ou refuser un devis, la mission n'etant creee qu'a
  l'acceptation ; affecter un vehicule et un chauffeur, au besoin avec
  l'assistance du moteur d'optimisation ; transmettre la position depuis
  l'application mobile ; suivre une mission en direct et consulter
  l'historique du trajet ; cloturer ou annuler une mission ; ouvrir et
  traiter un ticket de support \\

Maintenance curative et predictive &
  Ouvrir, instruire et valider un ordre de travail et en suivre
  l'avancement par etat ; declarer une intervention, la suivre, consommer
  ou restituer des pieces ; tenir le catalogue des pieces et
  reapprovisionner le stock avec alerte sur les seuils bas ; documenter
  un ordre de travail ; consulter le risque de panne par vehicule et
  l'organe concerne \\

Facturation, alertes et journal metier &
  Etablir une facture a partir des missions facturables, avec controle
  des doublons ; suivre le cycle de vie d'une facture et identifier les
  impayes ; consulter ses factures depuis le portail client ; recevoir
  les alertes metier portant sur les echeances, les stocks, les retards
  et les impayes ; consulter le journal des evenements et editer un
  rapport de synthese \\

Collaboration et aide a la decision &
  Echanger par messagerie instantanee, en conversation privee ou en
  groupe, avec pieces jointes ; consulter les correspondants connectes ;
  interroger un assistant documentaire sur les services et les donnees ;
  consulter un tableau de bord consolidant les indicateurs
  d'exploitation \\

\bottomrule
\end{tabularx}
\end{table}

Ces six domaines ne sont pas un simple classement de redaction :
\essentiel{ils determinent le decoupage en sprints} presente en
section~\ref{sec:sprints}, et chacun donnera son chapitre a la partie
consacree a la realisation.

% =====================================================================
\section{Besoins non fonctionnels}
\label{sec:besoins-non-fonctionnels}

Les besoins non fonctionnels ne decrivent pas ce que fait le systeme mais
la maniere dont il doit le faire. Ils ont oriente les choix
d'architecture presentes au chapitre suivant.

\begin{table}[H]
\centering
\caption{Besoins non fonctionnels et reponses apportees}
\label{tab:besoins-nf}
\small
\begin{tabularx}{\textwidth}{@{}p{2.6cm}X@{}}
\toprule
\textbf{Exigence} & \textbf{Contenu et reponse retenue} \\ \midrule
Securite & Jeton signe, second facteur optionnel, controle d'acces par
  role, verrouillage apres echecs repetes, journalisation des connexions.
  Aucune suppression definitive d'un compte porteur d'historique \\
Tracabilite & Toute operation sensible laisse une trace : historique des
  statuts d'ordre de travail, journal des connexions, journal metier.
  Reponse directe a la volatilite des echanges informels \\
Performance & Pagination systematique et filtrage cote serveur : les
  ecrans restent utilisables sur plusieurs milliers de lignes \\
Disponibilite & Chaque service expose un point de controle de sante ;
  un service defaillant n'empeche pas les autres de fonctionner \\
Maintenabilite & Responsabilites separees dans chaque service, services
  independants entre eux : faire evoluer la facturation n'oblige pas a
  redeployer la maintenance \\
Portabilite & Deploiement par conteneurs, sans installation manuelle. Le
  poste et le serveur executent la meme definition \\
Ergonomie & Trois interfaces plutot qu'une : chaque public ne voit que ce
  qui le concerne. Le chauffeur dispose d'une application mobile dediee \\
Evolutivite & Ajouter un domaine ne doit pas imposer de reprendre
  l'existant. Le huitieme service, la maintenance predictive, a ete
  ajoute apres coup : il valide ce point \\
\bottomrule
\end{tabularx}
\end{table}

% =====================================================================
\section{Modelisation des besoins : cas d'utilisation global}
\label{sec:uc-global}

Le diagramme de la figure~\ref{fig:uc-global} synthetise les besoins
fonctionnels. Il donne la carte d'ensemble de la plateforme ; le detail
de chaque domaine est presente dans le chapitre du sprint correspondant.

\figReelle{platform-usecase-global.png}
  {Diagramme de cas d'utilisation global de la plateforme \sujetCourt}
  {fig:uc-global}{0.60}

\paragraph{Lecture du diagramme.}
Les acteurs sont alignes a gauche, relies entre eux par des liens de
generalisation dont le triangle creux pointe vers l'acteur le plus
general. La frontiere du systeme est subdivisee en six domaines
fonctionnels, qui correspondent aux six sections des besoins
fonctionnels. Les deux systemes tiers figurent a droite, hors de la
frontiere.

\paragraph{Les cas d'aide a la decision.}
Trois cas sont volontairement places hors des domaines :
\exemple{estimer le prix du trajet}, \exemple{optimiser l'affectation de
la flotte} et \exemple{anticiper les pannes}. Ils ne constituent pas un
but d'utilisateur en eux-memes : personne ne se connecte pour
\exemple{lancer une prediction}. Ils sont atteints \essentiel{depuis un
cas de base}, au moment ou la decision se prend. C'est ce que traduisent
les relations d'inclusion et d'extension :

\begin{itemize}
  \item deposer une demande de transport \essentiel{inclut}
    systematiquement l'estimation du prix : la fourchette tarifaire est
    affichee au client avant qu'il ne valide ;
  \item etablir un devis est \essentiel{etendu} par l'estimation :
    le gestionnaire peut retenir le prix propose, le modifier ou le fixer
    lui-meme, l'estimation reste facultative ;
  \item planifier une mission \essentiel{inclut} l'optimisation, et
    piloter les ordres de travail \essentiel{inclut} l'anticipation des
    pannes.
\end{itemize}

\paragraph{Ce que le diagramme ne montre pas.}
Les services d'aide a la decision sont ici \essentiel{internes} au
systeme : ce sont des composants de la plateforme, pas des entites
exterieures. Seuls le serveur de messagerie et la passerelle de
messagerie instantanee, qui appartiennent a des tiers, figurent comme
acteurs secondaires.

% =====================================================================
\section{Pilotage du projet avec Scrum}
\label{sec:pilotage}

\subsection{Equipe et roles}

Le cadre Scrum a ete transpose au contexte d'un projet de fin d'etudes
selon la repartition du tableau~\ref{tab:equipe}.

\begin{table}[H]
\centering
\caption{Equipe projet et repartition des roles Scrum}
\label{tab:equipe}
\small
\begin{tabularx}{\textwidth}{@{}p{3.4cm}p{4.2cm}X@{}}
\toprule
\textbf{Role Scrum} & \textbf{Tenu par} & \textbf{Responsabilites} \\ \midrule
Product Owner & \encadrantEntreprise{} \newline (encadrant entreprise) &
  Porte le besoin metier, ordonne le carnet, valide les increments \\
Scrum Master & \encadrantAcademique{} \newline (encadrant academique) &
  Garantit le cadre et la rigueur methodologique, leve les obstacles \\
Equipe de developpement & \etudiantComplet &
  Conception, developpement, tests et deploiement \\
\bottomrule
\end{tabularx}
\end{table}

\subsection{Carnet de produit}

Le carnet de produit rassemble les besoins exprimes sous forme de recits
utilisateur, ordonnes par priorite. La priorisation suit la methode
MoSCoW : \exemple{Must have} pour ce qui est indispensable,
\exemple{Should have} pour ce qui est important sans etre bloquant,
\exemple{Could have} pour ce qui apporte un confort. L'estimation est
exprimee en points, selon une suite de Fibonacci.

\begin{table}[H]
\centering
\caption{Carnet de produit, par theme}
\label{tab:backlog-produit}
\small
\begin{tabularx}{\textwidth}{@{}p{0.9cm}Xp{1.9cm}p{1.4cm}@{}}
\toprule
\textbf{Ref.} & \textbf{Theme} & \textbf{Priorite} & \textbf{Points} \\
\midrule
T1 & Identites, acces et securite            & Must have   & 34 \\
T2 & Referentiel de flotte et d'equipements  & Must have   & 34 \\
T3 & Demande, devis, mission et suivi        & Must have   & 55 \\
T4 & Maintenance curative                    & Must have   & 34 \\
T5 & Facturation                             & Must have   & 21 \\
T6 & Alertes metier                          & Should have & 13 \\
T7 & Messagerie interne                      & Should have & 21 \\
T8 & Estimation tarifaire assistee           & Should have & 13 \\
T9 & Suivi de position en direct             & Should have & 21 \\
T10 & Maintenance predictive                 & Could have  & 13 \\
T11 & Journal metier et rapports             & Could have  & 13 \\
T12 & Assistant documentaire                 & Could have  & 21 \\
T13 & Optimisation de l'affectation          & Could have  & 13 \\
\midrule
\multicolumn{3}{@{}r}{\textbf{Total}} & \textbf{306} \\
\bottomrule
\end{tabularx}
\end{table}

\subsection{Decoupage en sprints}
\label{sec:sprints}

Le stage s'etendant sur \stageDuree{} mois, le travail a ete decoupe en
six sprints, precedes d'une phase de cadrage et suivis d'une phase de
recette. Chaque sprint couvre un \essentiel{domaine metier coherent} et
livre un increment utilisable de bout en bout : interface, service,
persistance.

\begin{table}[H]
\centering
\caption{Contenu des six sprints}
\label{tab:sprints}
\small
% La colonne « Duree » valait 1,2 cm : « 4 semaines » n'y tenait pas et
% se coupait en « 4 se-maines ». Portee a 2,3 cm, prise sur la colonne
% flexible qui a de la marge.
\begin{tabularx}{\textwidth}{@{}p{1.5cm}Xp{2.6cm}p{2.3cm}@{}}
\toprule
\textbf{Sprint} & \textbf{Domaine couvert} & \textbf{Themes} & \textbf{Duree} \\
\midrule
Sprint 1 & Identites, acces et securite & T1 & 4 semaines \\
Sprint 2 & Referentiel de flotte et d'equipements & T2 & 4 semaines \\
Sprint 3 & Demande, devis, mission, suivi en direct & T3, T8, T9 & 4 semaines \\
Sprint 4 & Maintenance curative et predictive & T4, T10 & 4 semaines \\
Sprint 5 & Facturation, alertes et journal metier & T5, T6, T11 & 3 semaines \\
Sprint 6 & Collaboration et aide a la decision & T7, T12, T13 & 3 semaines \\
\bottomrule
\end{tabularx}
\end{table}

L'ordre n'est pas arbitraire. \essentiel{Le sprint 1 conditionne tous les
autres} : sans identites ni habilitations, aucun ecran ne peut etre
protege. \essentiel{Le sprint 2 fournit le referentiel} sur lequel
s'appuient les missions et la maintenance : on ne peut affecter un
vehicule qui n'existe pas. Les sprints suivants exploitent ces
fondations, ce qui explique que les themes d'aide a la decision soient
traites au sein du sprint du domaine qu'ils servent, et non isolement.

\figReelle{planning-sprints.png}
  {Planification previsionnelle du projet}
  {fig:planning}{1.00}

Le diagramme de Gantt de la figure~\ref{fig:planning} situe ces phases
dans le calendrier du stage, du \stageDebut{} au \stageFin. La redaction
du rapport y apparait comme une activite transversale, engagee des le
quatrieme sprint plutot que reportee en fin de projet.

% =====================================================================
\section*{Conclusion}
\addcontentsline{toc}{section}{Conclusion}

Ce chapitre a traduit le besoin en exigences. Six acteurs humains
organises en hierarchie de generalisation et deux systemes tiers
interagissent avec la plateforme. Les besoins fonctionnels, presentes par
domaine, ont ete synthetises par un diagramme de cas d'utilisation global
dont la lecture met en evidence le parti pris principal : l'aide a la
decision n'est pas un module separe mais un ensemble de cas atteints
depuis les cas metier. Les besoins non fonctionnels ont ete formules et
rattaches aux reponses qui leur seront apportees. Le pilotage, enfin, a
ete arrete : trois roles Scrum, un carnet de produit de treize themes et
un decoupage en six sprints dont l'ordre suit les dependances entre
domaines.

Le chapitre suivant presente l'environnement de travail et l'architecture
qui repond a ces exigences.
